\section{The size estimate: shift lemma and certificates}\label{sec:size}

This section supplies the numerical input of (K3). A shift lemma moves a certificate to every smaller $z$
with a gain of $\frac12\log(z/z')$; four certificates, one for each of $n\in\{14,27\}$ and each sign,
are then computed, and a monotonicity argument extends them to all larger traces of the same parity.

\begin{lemma}[shift lemma]\label{lem:shift}
Let $0<z'\le z\le\sqrt5-2$ and $\Lambda:=\log(z/z')\ge0$. For every admissible $\mu$ and $\lambda\in\R$,
$\Phi_{z'}(\mu,\lambda-\Lambda)\le\Phi_z(\mu,\lambda)-\frac\Lambda2$. If $(\mu,\lambda)$ satisfies (i) at $z$ then
$(\mu,\lambda-\Lambda)$ satisfies (i) at $z'$; if it satisfies (ii) at $z$ with $c_{\mathrm{cell}}$, then
$(\mu,\lambda-\Lambda)$ satisfies (ii) at $z'$ with $c_{\mathrm{cell}}+1.01\Lambda$.
\end{lemma}

\begin{proof}
This is \cite[Lemma~5.5]{Krisman2026units}; the proof is short. Since $V_{z'}(x)=V_z(x)-\Lambda x$, we have
$F'_\mu-(\lambda-\Lambda)=(F_\mu-\lambda)-\Lambda(x-1)$. Put $a=F_\mu(x)-\lambda$. On $[0,1]$,
$(a+\Lambda(1-x))_+\le a_++\Lambda(1-x)$, and on $[1,\infty)$, $(a-\Lambda(x-1))_+\le a_+$. Integrating,
$\int_0^\infty(F'_\mu-\lambda+\Lambda)_+\le\int_0^\infty(F_\mu-\lambda)_++\frac\Lambda2$, and
$\Phi_{z'}(\mu,\lambda-\Lambda)\le-E(\mu)+\lambda-\Lambda+\int(F_\mu-\lambda)_++\frac\Lambda2=\Phi_z(\mu,\lambda)-\frac\Lambda2$.
On $[4,\infty)$ the term $-\Lambda(x-1)$ is negative, so (i) is preserved; on $[-\frac1{100},0]$ it is at most
$\Lambda(1+\frac1{100})$.
\end{proof}

For the certificates we write $t_\sigma(n)=\log(1/z_\sigma(n))=\log\frac{n+\sqrt{n^2+4\sigma}}2$.

\begin{proposition}[certificates]\label{prop:certificates}
For each $(n,\sigma)\in\{14,27\}\times\{\pm1\}$ there is an admissible step density $\mu$ (160 equal cells of
$[0,2]$ with heights $H_i/10^4$, $H_i\in\Z$) and a rational $\lambda$ such that, with $z=z_\sigma(n)\le\sqrt5-2$,
\textup{(i)} and \textup{(ii)} of \textup{(K3)} hold with $c_{\mathrm{cell}}\le7.41$, and
$\Car(n)+\Phi_z(\mu,\lambda)\le-0.005$. The values, rounded outward to six decimals, are in \cref{tab:certificates}; nine decimals are in \cref{app:certificates}.
\end{proposition}

\begin{table}[ht]
\centering\small
\caption{Certified comparisons, upper bounds rounded outward ($\Car(14)=1.993097\ldots$, $\Car(27)=2.330579\ldots$).}\label{tab:certificates}
\begin{tabular}{@{}rrrrrr@{}}
\toprule
$n$ & $\sigma$ & $\lambda$ & $\Phi_z(\mu,\lambda)\le$ & $\Car(n)+\Phi\le$ & $c_{\mathrm{cell}}\le$\\
\midrule
$14$ & $+1$ & $-4.3044$ & $-2.004611$ & $-0.011513$ & $6.73$\\
$14$ & $-1$ & $-4.2937$ & $-1.999368$ & $-0.006270$ & $6.72$\\
$27$ & $+1$ & $-4.9767$ & $-2.337342$ & $-0.006763$ & $7.41$\\
$27$ & $-1$ & $-4.9739$ & $-2.335955$ & $-0.005376$ & $7.41$\\
\bottomrule
\end{tabular}
\end{table}

\begin{proof}
A finite computation, described in \cref{app:certificates}: the densities and $\lambda$ are those of the supplementary
data files (see the section on reproducibility); the quantities $E(\mu)$, $U^\mu$ and the integral in \eqref{eq:Phi} are evaluated in
Arb ball arithmetic, the integral by an explicit upper bound on each mesh cell, and the side conditions as stated
there. The computation was done by two separate implementations (the one that produced the candidate and one written
for this paper); both give the same $\Car(n)+\Phi$ to $3.5\cdot10^{-7}$.
\end{proof}

\begin{proposition}\label{prop:all-traces}
Assume the data of \cref{prop:certificates}. Let $n_0=14$ and $n\ge14$ be even, or $n_0=27$ and $n\ge27$ be odd, with
the same $\sigma$. Put $z_0=z_\sigma(n_0)$, $z=z_\sigma(n)$, $\Lambda=t_\sigma(n)-t_\sigma(n_0)\ge0$ and let
$(\mu,\lambda)$ be the certificate of $(n_0,\sigma)$. Then $(\mu,\lambda-\Lambda)$ satisfies the hypotheses of \textup{(K3)}
at $z$ and
\[
\Car(n)+\Phi_z(\mu,\lambda-\Lambda)\le\Car(n_0)+\Phi_{z_0}(\mu,\lambda)<0 .
\]
\end{proposition}

\begin{proof}
Since $z_\sigma(\cdot)$ decreases, $0<z\le z_0\le\sqrt5-2$, and \cref{lem:shift} gives the hypotheses and
$\Phi_z(\mu,\lambda-\Lambda)\le\Phi_{z_0}(\mu,\lambda)-\frac\Lambda2$. By \cref{prop:arithmetic}, for even $n$,
$\Car(n)\le\frac43+\frac{\log n}4$, with equality at $n_0=14$; for odd $n$, $\Car(n)=\frac43+\frac{\log n}4+\frac{\log2}4$,
the same function of $n$ at $n_0=27$. So
$\Car(n)+\Phi_z\le\Car(n_0)+\Phi_{z_0}+[g(n)-g(n_0)]$ with $g(x)=\frac14\log x-\frac12t_\sigma(x)$,
$t_\sigma(x)=\log\frac{x+\sqrt{x^2+4\sigma}}2$. As $t_\sigma'(x)=1/\sqrt{x^2+4\sigma}$,
$g'(x)=\frac1{4x}-\frac1{2\sqrt{x^2+4\sigma}}<0$ for $x\ge2$, because $\sqrt{x^2+4\sigma}<2x$. Hence
$g(n)\le g(n_0)$.
\end{proof}
